% thesis-tablestyle.tex -- one table style for the whole thesis (every chapter, generated or hand-written).
%
% Load ONCE in the preamble, after xcolor (with its `table' option) and booktabs:
%     \usepackage[table,usenames,dvipsnames]{xcolor}   % already in dissertation.tex
%     \usepackage{booktabs}                            % already in dissertation.tex
%     % thesis-tablestyle.tex -- one table style for the whole thesis (every chapter, generated or hand-written).
%
% Load ONCE in the preamble, after xcolor (with its `table' option) and booktabs:
%     \usepackage[table,usenames,dvipsnames]{xcolor}   % already in dissertation.tex
%     \usepackage{booktabs}                            % already in dissertation.tex
%     % thesis-tablestyle.tex -- one table style for the whole thesis (every chapter, generated or hand-written).
%
% Load ONCE in the preamble, after xcolor (with its `table' option) and booktabs:
%     \usepackage[table,usenames,dvipsnames]{xcolor}   % already in dissertation.tex
%     \usepackage{booktabs}                            % already in dissertation.tex
%     % thesis-tablestyle.tex -- one table style for the whole thesis (every chapter, generated or hand-written).
%
% Load ONCE in the preamble, after xcolor (with its `table' option) and booktabs:
%     \usepackage[table,usenames,dvipsnames]{xcolor}   % already in dissertation.tex
%     \usepackage{booktabs}                            % already in dissertation.tex
%     \input{thesis-tablestyle}
%
% What it does, in every table / table* / sidewaystable / sidewaystable* float:
%   * zebra striping of the body rows (header rows stay white), in a very light grey that prints in greyscale;
%   * slightly taller rows (\arraystretch times \thesisrowstretch, on top of any \arraystretch the table sets);
%   * booktabs rules keep their spacing; where a rule follows a shaded row, the gap above the rule is shaded
%     too, so there is no white band.
% The body starts after the first \midrule (booktabs) or the second \hline (\hline tables). Striping restarts
% after every \midrule, so a group row after a \midrule (e.g. "\multicolumn{8}{l}{\textit{PropNEAT}}") is
% white and its rows alternate below it. A table whose header ends without a \midrule says how many header
% rows it has: \thesistablestyle[1]. A table with no top rule (\toprule or \hline) is not striped.
%
% Switches (in the preamble for the whole thesis, or inside one table for that table):
%   \thesisstripestrue / \thesisstripesfalse        zebra on (default) / off
%   \thesisfiverulestrue / \thesisfiverulesfalse    thin grey rule every \thesisfiveevery body rows (default off);
%                                                   the alternative to stripes: \thesisstripesfalse\thesisfiverulestrue
%   \thesistablesautotrue / \thesistablesautofalse  apply to every table float automatically (default on); when
%                                                   off, only tables that call \thesistablestyle are styled
%   \renewcommand\thesisrowstretch{1.1}             row-height factor (1 = unchanged)
%   \renewcommand\thesisfiveevery{5}                rows between the thin rules
%   \definecolor{thesisstripe}{gray}{0.93}          stripe colour (and thesisfiverule for the thin rules)
% Per table:
%   \thesistablestyle[<n>]   apply the style (generated tables call it first; idempotent); <n> = number of
%                            header rows after \toprule when the header does not end with a \midrule
%   \thesisnostripes         opt out: this table keeps its own look (no stripes, no thin rules, no extra
%                            height). Use it for tables that colour their own cells or rows, or use \multirow
%                            (the next row's shading would cover the \multirow text).
% Nested tabulars (\makecell, a tabular inside a cell) are left alone and do not disturb the outer striping.
\makeatletter
\@ifpackageloaded{colortbl}{}{\RequirePackage{colortbl}}
\@ifpackageloaded{booktabs}{}{\RequirePackage{booktabs}}
\@ifundefined{fpeval}{\RequirePackage{xfp}}{}

\newif\ifthesisstripes   \thesisstripestrue
\newif\ifthesisfiverules \thesisfiverulesfalse
\newif\ifthesistablesauto \thesistablesautotrue
\definecolor{thesisstripe}{gray}{0.93}
\definecolor{thesisfiverule}{gray}{0.6}
\newcommand\thesisrowstretch{1.1}
\newcommand\thesisfiveevery{5}

\newif\ifthesis@active   % inside a styled float
\newif\ifthesis@engine   % stripes or thin rules are being drawn
\newif\ifthesis@data     % past the header
\newif\ifthesis@firstrule
\newif\ifthesis@lastrow  % the last item was a row (not a rule)
\newif\ifthesis@curgrey  % the current row is shaded
\newif\ifthesis@lastgrey % the row just ended was shaded
\newif\ifthesis@wantsep
\newif\ifthesis@oksep
\newcount\thesis@row     % body row number since the last \midrule (1 = first)
\newcount\thesis@five
\newcount\thesis@head    % header rows still to come (-1: wait for a \midrule)
\newcount\thesis@lvl     % group level of the outer table's \noalign (-1: unknown)
\thesis@lvl=\m@ne
\def\thesis@headrows{-1}

% ---- row colours (colortbl's per-row colour, as xcolor's \rowcolors sets it). colortbl extends each
% cell's colour by \tabcolsep on both sides; at an @{} edge that would stick out past the rules, so the
% first / last cell of a row drops that overhang when the column spec starts / ends with @{}.
\newif\ifthesis@leftflush
\newif\ifthesis@rightflush
\newcount\thesis@cellno
\newcount\thesis@ncols
% colortbl stretches a cell's colour with 1fill against the cell's own 1fill of alignment glue; a
% \makecell adds its own \hfill glue (counted in \thesis@mc), which would take part of the stretch and
% leave the right of the cell white. The colour's stretch is raised to match (\thesis@k).
\newcount\thesis@mc
\def\thesis@k{1}
\def\thesis@mcell#1#2\@nil#3{\thesis@o@mcell{#1}#2\@nil{#3}%
  \ifx\mcell@left\hfill\global\advance\thesis@mc\@ne\fi
  \ifx\mcell@right\hfill\global\advance\thesis@mc\@ne\fi}
\def\thesis@docolor{\global\let\CT@do@color\relax
  \@tempdima\wd\z@\advance\@tempdima\@tempdimb\advance\@tempdima\@tempdimc
  \kern-\@tempdimb\leaders\vrule\hskip\@tempdima\@plus\thesis@k fill
  \kern-\@tempdimc\hskip-\wd\z@\@plus-\thesis@k fill\relax}
\def\thesis@cellstep{%
  \def\thesis@k{1}%
  \ifnum\currentgrouplevel=\thesis@lvl
    \edef\thesis@k{\the\numexpr\@ne+\thesis@mc\relax}\global\thesis@mc\z@
    \global\advance\thesis@cellno\@ne
    % overlap the next cell a little, so screen renderers show no hairline gaps between cells
    \ifnum\thesis@cellno<\thesis@ncols\advance\@tempdimc.25pt\fi
    \ifnum\thesis@cellno=\@ne\ifthesis@leftflush\@tempdimb\z@\fi\fi
    \ifnum\thesis@cellno=\thesis@ncols\ifthesis@rightflush\@tempdimc\z@\fi\fi
  \fi}
\def\thesis@grey{\gdef\CT@row@color{\thesis@cellstep\CT@color{thesisstripe}}\global\thesis@curgreytrue}
\def\thesis@white{\gdef\CT@row@color{\thesis@cellstep}\global\thesis@curgreyfalse}
\ExplSyntaxOn
\cs_new_protected:Npn \thesis@edges #1
  {
    \str_set:Nn \l_tmpa_str {#1}
    \str_remove_all:Nn \l_tmpa_str { ~ }
    \str_if_eq:eeTF { \str_head:N \l_tmpa_str } { @ }
      { \global\thesis@leftflushtrue } { \global\thesis@leftflushfalse }
    \str_if_eq:eeTF { \str_range:Nnn \l_tmpa_str { -3 } { -1 } } { \tl_to_str:n { @{} } }
      { \global\thesis@rightflushtrue } { \global\thesis@rightflushfalse }
  }
\ExplSyntaxOff
\def\thesis@array[#1]#2{%
  \ifnum\thesis@lvl<\z@\thesis@edges{#2}\else\ifnum\currentgrouplevel<\thesis@lvl\thesis@edges{#2}\fi\fi
  \thesis@o@array[#1]{#2}}
\def\thesis@colour{\ifthesisstripes\ifodd\thesis@row\thesis@white\else\thesis@grey\fi\else\thesis@white\fi}
\def\thesis@startdata{\global\thesis@datatrue\global\thesis@head\m@ne\global\thesis@row\@ne
  \global\thesis@five\z@\thesis@colour}
\def\thesis@sethead{\global\thesis@lvl\currentgrouplevel\global\thesis@firstrulefalse
  \global\thesis@ncols\z@\global\thesis@cellno\z@
  \global\thesis@datafalse\global\thesis@row\z@\global\thesis@five\z@\thesis@white
  \global\thesis@head=\thesis@headrows\relax\ifnum\thesis@head=\z@\thesis@startdata\fi}

% ---- at every row end (\CT@everycr); only the outer table's rows count
\def\thesis@close{\ifnum0=`{\fi}}
\def\thesis@cr{\noalign{\ifnum0=`}\fi\thesis@crstep}
\def\thesis@crstep{\thesis@wantsepfalse
  \ifnum\currentgrouplevel=\thesis@lvl
    \ifthesis@curgrey\global\thesis@lastgreytrue\else\global\thesis@lastgreyfalse\fi
    \global\thesis@lastrowtrue
    \ifnum\thesis@cellno>\thesis@ncols\global\thesis@ncols\thesis@cellno\fi
    \global\thesis@cellno\z@\global\thesis@mc\z@
    \ifthesis@data
      \global\advance\thesis@row\@ne\thesis@colour
      \ifthesisfiverules
        \global\advance\thesis@five\@ne
        \ifnum\thesis@five<\thesisfiveevery\relax\else\global\thesis@five\z@\thesis@wantseptrue\fi
      \fi
    \else
      \thesis@white
      \ifnum\thesis@head>\z@
        \global\advance\thesis@head\m@ne
        \ifnum\thesis@head=\z@\thesis@startdata\fi
      \fi
    \fi
  \fi
  \ifthesis@wantsep\expandafter\thesis@peek\else\expandafter\thesis@close\fi}
% a thin rule after every \thesisfiveevery body rows, unless a rule or the end comes next
\def\thesis@peek{\futurenonspacelet\thesis@next\thesis@decide}
\def\thesis@nosepif#1{\ifx\thesis@next#1\thesis@oksepfalse\fi}
\def\thesis@decide{\thesis@okseptrue
  \thesis@nosepif\bottomrule\thesis@nosepif\midrule\thesis@nosepif\toprule\thesis@nosepif\cmidrule
  \thesis@nosepif\hline\thesis@nosepif\cline\thesis@nosepif\addlinespace\thesis@nosepif\specialrule
  \thesis@nosepif\end\thesis@nosepif\noalign\thesis@nosepif\crcr
  \ifthesis@oksep{\color{thesisfiverule}\hrule\@height\cmidrulewidth}\fi
  \thesis@close}

% ---- rules: #1 = 0 \cmidrule, 1 \toprule, 2 \midrule, 3 \bottomrule, 4 \hline
\def\thesis@fill{{\color{thesisstripe}\hrule\@height\aboverulesep}\vskip-\aboverulesep}
\def\thesis@rule#1{\noalign{\thesis@rulestep#1}}
\def\thesis@rulestep#1{%
  \ifnum#1=\@ne\thesis@sethead\else
  \ifnum#1=4
    \ifthesis@firstrule\thesis@sethead\else
    \ifnum\currentgrouplevel=\thesis@lvl\thesis@startdata\else\thesis@sethead\fi\fi
  \else
  \ifnum\currentgrouplevel=\thesis@lvl
    \ifthesis@lastrow\ifthesis@lastgrey\thesis@fill\fi\fi
    \ifnum#1=\tw@\thesis@startdata\fi
    \ifnum#1=\thr@@\global\thesis@datafalse\thesis@white\fi
  \fi\fi\fi
  \global\thesis@lastrowfalse}

% ---- at the start of each tabular / tabularx / tabular* (not nested ones): reset, taller rows
\def\thesis@tabbegin{%
  \ifthesis@active
    \ifnum\thesis@lvl<\z@\thesis@tabreset\else
    \ifnum\currentgrouplevel<\thesis@lvl\thesis@tabreset\fi\fi
  \fi}
\def\thesis@tabreset{\global\thesis@lvl\m@ne\global\thesis@firstruletrue\global\thesis@datafalse
  \global\thesis@row\z@\global\thesis@five\z@\global\thesis@head\m@ne\global\thesis@lastrowfalse
  \global\thesis@curgreyfalse
  \edef\arraystretch{\fpeval{\arraystretch*\thesisrowstretch}}}
\AddToHook{env/tabular/begin}{\thesis@tabbegin}
\AddToHook{env/tabular*/begin}{\thesis@tabbegin}
\AddToHook{env/tabularx/begin}{\thesis@tabbegin}

% ---- the style itself
\newcommand\thesistablestyle[1][-1]{%
  \def\thesis@headrows{#1}%
  \ifthesis@active\else
    \thesis@activetrue\global\thesis@lvl\m@ne
    \thesis@enginefalse
    \ifthesisstripes\thesis@enginetrue\fi
    \ifthesisfiverules\thesis@enginetrue\fi
    \ifthesis@engine
      \CT@everycr{\the\everycr\thesis@cr}%
      \let\thesis@o@toprule\toprule\def\toprule{\thesis@rule1\thesis@o@toprule}%
      \let\thesis@o@midrule\midrule\def\midrule{\thesis@rule2\thesis@o@midrule}%
      \let\thesis@o@bottomrule\bottomrule\def\bottomrule{\thesis@rule3\thesis@o@bottomrule}%
      \let\thesis@o@cmidrule\cmidrule\def\cmidrule{\thesis@rule0\thesis@o@cmidrule}%
      \let\thesis@o@hline\hline\def\hline{\thesis@rule4\thesis@o@hline}%
      \ifdefined\@@array\let\thesis@o@array\@@array\let\@@array\thesis@array\fi
      \ifdefined\mcell@@@tabular\let\thesis@o@mcell\mcell@@@tabular\let\mcell@@@tabular\thesis@mcell\fi
      \let\thesis@o@docolor\CT@@do@color\let\CT@@do@color\thesis@docolor
    \fi
  \fi}
\newcommand\thesisnostripes{%
  \ifthesis@active
    \thesis@activefalse
    \ifthesis@engine
      \CT@everycr{\noalign{\global\let\CT@row@color\relax}\the\everycr}%
      \let\toprule\thesis@o@toprule\let\midrule\thesis@o@midrule\let\bottomrule\thesis@o@bottomrule
      \let\cmidrule\thesis@o@cmidrule\let\hline\thesis@o@hline
      \ifdefined\thesis@o@array\let\@@array\thesis@o@array\fi
      \ifdefined\thesis@o@mcell\let\mcell@@@tabular\thesis@o@mcell\fi
      \let\CT@@do@color\thesis@o@docolor
      \thesis@enginefalse
    \fi
  \fi}
\def\thesis@auto{\ifthesistablesauto\thesistablestyle\fi}
\AddToHook{env/table/begin}{\thesis@auto}
\AddToHook{env/table*/begin}{\thesis@auto}
\AddToHook{env/sidewaystable/begin}{\thesis@auto}
\AddToHook{env/sidewaystable*/begin}{\thesis@auto}
\makeatother

%
% What it does, in every table / table* / sidewaystable / sidewaystable* float:
%   * zebra striping of the body rows (header rows stay white), in a very light grey that prints in greyscale;
%   * slightly taller rows (\arraystretch times \thesisrowstretch, on top of any \arraystretch the table sets);
%   * booktabs rules keep their spacing; where a rule follows a shaded row, the gap above the rule is shaded
%     too, so there is no white band.
% The body starts after the first \midrule (booktabs) or the second \hline (\hline tables). Striping restarts
% after every \midrule, so a group row after a \midrule (e.g. "\multicolumn{8}{l}{\textit{PropNEAT}}") is
% white and its rows alternate below it. A table whose header ends without a \midrule says how many header
% rows it has: \thesistablestyle[1]. A table with no top rule (\toprule or \hline) is not striped.
%
% Switches (in the preamble for the whole thesis, or inside one table for that table):
%   \thesisstripestrue / \thesisstripesfalse        zebra on (default) / off
%   \thesisfiverulestrue / \thesisfiverulesfalse    thin grey rule every \thesisfiveevery body rows (default off);
%                                                   the alternative to stripes: \thesisstripesfalse\thesisfiverulestrue
%   \thesistablesautotrue / \thesistablesautofalse  apply to every table float automatically (default on); when
%                                                   off, only tables that call \thesistablestyle are styled
%   \renewcommand\thesisrowstretch{1.1}             row-height factor (1 = unchanged)
%   \renewcommand\thesisfiveevery{5}                rows between the thin rules
%   \definecolor{thesisstripe}{gray}{0.93}          stripe colour (and thesisfiverule for the thin rules)
% Per table:
%   \thesistablestyle[<n>]   apply the style (generated tables call it first; idempotent); <n> = number of
%                            header rows after \toprule when the header does not end with a \midrule
%   \thesisnostripes         opt out: this table keeps its own look (no stripes, no thin rules, no extra
%                            height). Use it for tables that colour their own cells or rows, or use \multirow
%                            (the next row's shading would cover the \multirow text).
% Nested tabulars (\makecell, a tabular inside a cell) are left alone and do not disturb the outer striping.
\makeatletter
\@ifpackageloaded{colortbl}{}{\RequirePackage{colortbl}}
\@ifpackageloaded{booktabs}{}{\RequirePackage{booktabs}}
\@ifundefined{fpeval}{\RequirePackage{xfp}}{}

\newif\ifthesisstripes   \thesisstripestrue
\newif\ifthesisfiverules \thesisfiverulesfalse
\newif\ifthesistablesauto \thesistablesautotrue
\definecolor{thesisstripe}{gray}{0.93}
\definecolor{thesisfiverule}{gray}{0.6}
\newcommand\thesisrowstretch{1.1}
\newcommand\thesisfiveevery{5}

\newif\ifthesis@active   % inside a styled float
\newif\ifthesis@engine   % stripes or thin rules are being drawn
\newif\ifthesis@data     % past the header
\newif\ifthesis@firstrule
\newif\ifthesis@lastrow  % the last item was a row (not a rule)
\newif\ifthesis@curgrey  % the current row is shaded
\newif\ifthesis@lastgrey % the row just ended was shaded
\newif\ifthesis@wantsep
\newif\ifthesis@oksep
\newcount\thesis@row     % body row number since the last \midrule (1 = first)
\newcount\thesis@five
\newcount\thesis@head    % header rows still to come (-1: wait for a \midrule)
\newcount\thesis@lvl     % group level of the outer table's \noalign (-1: unknown)
\thesis@lvl=\m@ne
\def\thesis@headrows{-1}

% ---- row colours (colortbl's per-row colour, as xcolor's \rowcolors sets it). colortbl extends each
% cell's colour by \tabcolsep on both sides; at an @{} edge that would stick out past the rules, so the
% first / last cell of a row drops that overhang when the column spec starts / ends with @{}.
\newif\ifthesis@leftflush
\newif\ifthesis@rightflush
\newcount\thesis@cellno
\newcount\thesis@ncols
% colortbl stretches a cell's colour with 1fill against the cell's own 1fill of alignment glue; a
% \makecell adds its own \hfill glue (counted in \thesis@mc), which would take part of the stretch and
% leave the right of the cell white. The colour's stretch is raised to match (\thesis@k).
\newcount\thesis@mc
\def\thesis@k{1}
\def\thesis@mcell#1#2\@nil#3{\thesis@o@mcell{#1}#2\@nil{#3}%
  \ifx\mcell@left\hfill\global\advance\thesis@mc\@ne\fi
  \ifx\mcell@right\hfill\global\advance\thesis@mc\@ne\fi}
\def\thesis@docolor{\global\let\CT@do@color\relax
  \@tempdima\wd\z@\advance\@tempdima\@tempdimb\advance\@tempdima\@tempdimc
  \kern-\@tempdimb\leaders\vrule\hskip\@tempdima\@plus\thesis@k fill
  \kern-\@tempdimc\hskip-\wd\z@\@plus-\thesis@k fill\relax}
\def\thesis@cellstep{%
  \def\thesis@k{1}%
  \ifnum\currentgrouplevel=\thesis@lvl
    \edef\thesis@k{\the\numexpr\@ne+\thesis@mc\relax}\global\thesis@mc\z@
    \global\advance\thesis@cellno\@ne
    % overlap the next cell a little, so screen renderers show no hairline gaps between cells
    \ifnum\thesis@cellno<\thesis@ncols\advance\@tempdimc.25pt\fi
    \ifnum\thesis@cellno=\@ne\ifthesis@leftflush\@tempdimb\z@\fi\fi
    \ifnum\thesis@cellno=\thesis@ncols\ifthesis@rightflush\@tempdimc\z@\fi\fi
  \fi}
\def\thesis@grey{\gdef\CT@row@color{\thesis@cellstep\CT@color{thesisstripe}}\global\thesis@curgreytrue}
\def\thesis@white{\gdef\CT@row@color{\thesis@cellstep}\global\thesis@curgreyfalse}
\ExplSyntaxOn
\cs_new_protected:Npn \thesis@edges #1
  {
    \str_set:Nn \l_tmpa_str {#1}
    \str_remove_all:Nn \l_tmpa_str { ~ }
    \str_if_eq:eeTF { \str_head:N \l_tmpa_str } { @ }
      { \global\thesis@leftflushtrue } { \global\thesis@leftflushfalse }
    \str_if_eq:eeTF { \str_range:Nnn \l_tmpa_str { -3 } { -1 } } { \tl_to_str:n { @{} } }
      { \global\thesis@rightflushtrue } { \global\thesis@rightflushfalse }
  }
\ExplSyntaxOff
\def\thesis@array[#1]#2{%
  \ifnum\thesis@lvl<\z@\thesis@edges{#2}\else\ifnum\currentgrouplevel<\thesis@lvl\thesis@edges{#2}\fi\fi
  \thesis@o@array[#1]{#2}}
\def\thesis@colour{\ifthesisstripes\ifodd\thesis@row\thesis@white\else\thesis@grey\fi\else\thesis@white\fi}
\def\thesis@startdata{\global\thesis@datatrue\global\thesis@head\m@ne\global\thesis@row\@ne
  \global\thesis@five\z@\thesis@colour}
\def\thesis@sethead{\global\thesis@lvl\currentgrouplevel\global\thesis@firstrulefalse
  \global\thesis@ncols\z@\global\thesis@cellno\z@
  \global\thesis@datafalse\global\thesis@row\z@\global\thesis@five\z@\thesis@white
  \global\thesis@head=\thesis@headrows\relax\ifnum\thesis@head=\z@\thesis@startdata\fi}

% ---- at every row end (\CT@everycr); only the outer table's rows count
\def\thesis@close{\ifnum0=`{\fi}}
\def\thesis@cr{\noalign{\ifnum0=`}\fi\thesis@crstep}
\def\thesis@crstep{\thesis@wantsepfalse
  \ifnum\currentgrouplevel=\thesis@lvl
    \ifthesis@curgrey\global\thesis@lastgreytrue\else\global\thesis@lastgreyfalse\fi
    \global\thesis@lastrowtrue
    \ifnum\thesis@cellno>\thesis@ncols\global\thesis@ncols\thesis@cellno\fi
    \global\thesis@cellno\z@\global\thesis@mc\z@
    \ifthesis@data
      \global\advance\thesis@row\@ne\thesis@colour
      \ifthesisfiverules
        \global\advance\thesis@five\@ne
        \ifnum\thesis@five<\thesisfiveevery\relax\else\global\thesis@five\z@\thesis@wantseptrue\fi
      \fi
    \else
      \thesis@white
      \ifnum\thesis@head>\z@
        \global\advance\thesis@head\m@ne
        \ifnum\thesis@head=\z@\thesis@startdata\fi
      \fi
    \fi
  \fi
  \ifthesis@wantsep\expandafter\thesis@peek\else\expandafter\thesis@close\fi}
% a thin rule after every \thesisfiveevery body rows, unless a rule or the end comes next
\def\thesis@peek{\futurenonspacelet\thesis@next\thesis@decide}
\def\thesis@nosepif#1{\ifx\thesis@next#1\thesis@oksepfalse\fi}
\def\thesis@decide{\thesis@okseptrue
  \thesis@nosepif\bottomrule\thesis@nosepif\midrule\thesis@nosepif\toprule\thesis@nosepif\cmidrule
  \thesis@nosepif\hline\thesis@nosepif\cline\thesis@nosepif\addlinespace\thesis@nosepif\specialrule
  \thesis@nosepif\end\thesis@nosepif\noalign\thesis@nosepif\crcr
  \ifthesis@oksep{\color{thesisfiverule}\hrule\@height\cmidrulewidth}\fi
  \thesis@close}

% ---- rules: #1 = 0 \cmidrule, 1 \toprule, 2 \midrule, 3 \bottomrule, 4 \hline
\def\thesis@fill{{\color{thesisstripe}\hrule\@height\aboverulesep}\vskip-\aboverulesep}
\def\thesis@rule#1{\noalign{\thesis@rulestep#1}}
\def\thesis@rulestep#1{%
  \ifnum#1=\@ne\thesis@sethead\else
  \ifnum#1=4
    \ifthesis@firstrule\thesis@sethead\else
    \ifnum\currentgrouplevel=\thesis@lvl\thesis@startdata\else\thesis@sethead\fi\fi
  \else
  \ifnum\currentgrouplevel=\thesis@lvl
    \ifthesis@lastrow\ifthesis@lastgrey\thesis@fill\fi\fi
    \ifnum#1=\tw@\thesis@startdata\fi
    \ifnum#1=\thr@@\global\thesis@datafalse\thesis@white\fi
  \fi\fi\fi
  \global\thesis@lastrowfalse}

% ---- at the start of each tabular / tabularx / tabular* (not nested ones): reset, taller rows
\def\thesis@tabbegin{%
  \ifthesis@active
    \ifnum\thesis@lvl<\z@\thesis@tabreset\else
    \ifnum\currentgrouplevel<\thesis@lvl\thesis@tabreset\fi\fi
  \fi}
\def\thesis@tabreset{\global\thesis@lvl\m@ne\global\thesis@firstruletrue\global\thesis@datafalse
  \global\thesis@row\z@\global\thesis@five\z@\global\thesis@head\m@ne\global\thesis@lastrowfalse
  \global\thesis@curgreyfalse
  \edef\arraystretch{\fpeval{\arraystretch*\thesisrowstretch}}}
\AddToHook{env/tabular/begin}{\thesis@tabbegin}
\AddToHook{env/tabular*/begin}{\thesis@tabbegin}
\AddToHook{env/tabularx/begin}{\thesis@tabbegin}

% ---- the style itself
\newcommand\thesistablestyle[1][-1]{%
  \def\thesis@headrows{#1}%
  \ifthesis@active\else
    \thesis@activetrue\global\thesis@lvl\m@ne
    \thesis@enginefalse
    \ifthesisstripes\thesis@enginetrue\fi
    \ifthesisfiverules\thesis@enginetrue\fi
    \ifthesis@engine
      \CT@everycr{\the\everycr\thesis@cr}%
      \let\thesis@o@toprule\toprule\def\toprule{\thesis@rule1\thesis@o@toprule}%
      \let\thesis@o@midrule\midrule\def\midrule{\thesis@rule2\thesis@o@midrule}%
      \let\thesis@o@bottomrule\bottomrule\def\bottomrule{\thesis@rule3\thesis@o@bottomrule}%
      \let\thesis@o@cmidrule\cmidrule\def\cmidrule{\thesis@rule0\thesis@o@cmidrule}%
      \let\thesis@o@hline\hline\def\hline{\thesis@rule4\thesis@o@hline}%
      \ifdefined\@@array\let\thesis@o@array\@@array\let\@@array\thesis@array\fi
      \ifdefined\mcell@@@tabular\let\thesis@o@mcell\mcell@@@tabular\let\mcell@@@tabular\thesis@mcell\fi
      \let\thesis@o@docolor\CT@@do@color\let\CT@@do@color\thesis@docolor
    \fi
  \fi}
\newcommand\thesisnostripes{%
  \ifthesis@active
    \thesis@activefalse
    \ifthesis@engine
      \CT@everycr{\noalign{\global\let\CT@row@color\relax}\the\everycr}%
      \let\toprule\thesis@o@toprule\let\midrule\thesis@o@midrule\let\bottomrule\thesis@o@bottomrule
      \let\cmidrule\thesis@o@cmidrule\let\hline\thesis@o@hline
      \ifdefined\thesis@o@array\let\@@array\thesis@o@array\fi
      \ifdefined\thesis@o@mcell\let\mcell@@@tabular\thesis@o@mcell\fi
      \let\CT@@do@color\thesis@o@docolor
      \thesis@enginefalse
    \fi
  \fi}
\def\thesis@auto{\ifthesistablesauto\thesistablestyle\fi}
\AddToHook{env/table/begin}{\thesis@auto}
\AddToHook{env/table*/begin}{\thesis@auto}
\AddToHook{env/sidewaystable/begin}{\thesis@auto}
\AddToHook{env/sidewaystable*/begin}{\thesis@auto}
\makeatother

%
% What it does, in every table / table* / sidewaystable / sidewaystable* float:
%   * zebra striping of the body rows (header rows stay white), in a very light grey that prints in greyscale;
%   * slightly taller rows (\arraystretch times \thesisrowstretch, on top of any \arraystretch the table sets);
%   * booktabs rules keep their spacing; where a rule follows a shaded row, the gap above the rule is shaded
%     too, so there is no white band.
% The body starts after the first \midrule (booktabs) or the second \hline (\hline tables). Striping restarts
% after every \midrule, so a group row after a \midrule (e.g. "\multicolumn{8}{l}{\textit{PropNEAT}}") is
% white and its rows alternate below it. A table whose header ends without a \midrule says how many header
% rows it has: \thesistablestyle[1]. A table with no top rule (\toprule or \hline) is not striped.
%
% Switches (in the preamble for the whole thesis, or inside one table for that table):
%   \thesisstripestrue / \thesisstripesfalse        zebra on (default) / off
%   \thesisfiverulestrue / \thesisfiverulesfalse    thin grey rule every \thesisfiveevery body rows (default off);
%                                                   the alternative to stripes: \thesisstripesfalse\thesisfiverulestrue
%   \thesistablesautotrue / \thesistablesautofalse  apply to every table float automatically (default on); when
%                                                   off, only tables that call \thesistablestyle are styled
%   \renewcommand\thesisrowstretch{1.1}             row-height factor (1 = unchanged)
%   \renewcommand\thesisfiveevery{5}                rows between the thin rules
%   \definecolor{thesisstripe}{gray}{0.93}          stripe colour (and thesisfiverule for the thin rules)
% Per table:
%   \thesistablestyle[<n>]   apply the style (generated tables call it first; idempotent); <n> = number of
%                            header rows after \toprule when the header does not end with a \midrule
%   \thesisnostripes         opt out: this table keeps its own look (no stripes, no thin rules, no extra
%                            height). Use it for tables that colour their own cells or rows, or use \multirow
%                            (the next row's shading would cover the \multirow text).
% Nested tabulars (\makecell, a tabular inside a cell) are left alone and do not disturb the outer striping.
\makeatletter
\@ifpackageloaded{colortbl}{}{\RequirePackage{colortbl}}
\@ifpackageloaded{booktabs}{}{\RequirePackage{booktabs}}
\@ifundefined{fpeval}{\RequirePackage{xfp}}{}

\newif\ifthesisstripes   \thesisstripestrue
\newif\ifthesisfiverules \thesisfiverulesfalse
\newif\ifthesistablesauto \thesistablesautotrue
\definecolor{thesisstripe}{gray}{0.93}
\definecolor{thesisfiverule}{gray}{0.6}
\newcommand\thesisrowstretch{1.1}
\newcommand\thesisfiveevery{5}

\newif\ifthesis@active   % inside a styled float
\newif\ifthesis@engine   % stripes or thin rules are being drawn
\newif\ifthesis@data     % past the header
\newif\ifthesis@firstrule
\newif\ifthesis@lastrow  % the last item was a row (not a rule)
\newif\ifthesis@curgrey  % the current row is shaded
\newif\ifthesis@lastgrey % the row just ended was shaded
\newif\ifthesis@wantsep
\newif\ifthesis@oksep
\newcount\thesis@row     % body row number since the last \midrule (1 = first)
\newcount\thesis@five
\newcount\thesis@head    % header rows still to come (-1: wait for a \midrule)
\newcount\thesis@lvl     % group level of the outer table's \noalign (-1: unknown)
\thesis@lvl=\m@ne
\def\thesis@headrows{-1}

% ---- row colours (colortbl's per-row colour, as xcolor's \rowcolors sets it). colortbl extends each
% cell's colour by \tabcolsep on both sides; at an @{} edge that would stick out past the rules, so the
% first / last cell of a row drops that overhang when the column spec starts / ends with @{}.
\newif\ifthesis@leftflush
\newif\ifthesis@rightflush
\newcount\thesis@cellno
\newcount\thesis@ncols
% colortbl stretches a cell's colour with 1fill against the cell's own 1fill of alignment glue; a
% \makecell adds its own \hfill glue (counted in \thesis@mc), which would take part of the stretch and
% leave the right of the cell white. The colour's stretch is raised to match (\thesis@k).
\newcount\thesis@mc
\def\thesis@k{1}
\def\thesis@mcell#1#2\@nil#3{\thesis@o@mcell{#1}#2\@nil{#3}%
  \ifx\mcell@left\hfill\global\advance\thesis@mc\@ne\fi
  \ifx\mcell@right\hfill\global\advance\thesis@mc\@ne\fi}
\def\thesis@docolor{\global\let\CT@do@color\relax
  \@tempdima\wd\z@\advance\@tempdima\@tempdimb\advance\@tempdima\@tempdimc
  \kern-\@tempdimb\leaders\vrule\hskip\@tempdima\@plus\thesis@k fill
  \kern-\@tempdimc\hskip-\wd\z@\@plus-\thesis@k fill\relax}
\def\thesis@cellstep{%
  \def\thesis@k{1}%
  \ifnum\currentgrouplevel=\thesis@lvl
    \edef\thesis@k{\the\numexpr\@ne+\thesis@mc\relax}\global\thesis@mc\z@
    \global\advance\thesis@cellno\@ne
    % overlap the next cell a little, so screen renderers show no hairline gaps between cells
    \ifnum\thesis@cellno<\thesis@ncols\advance\@tempdimc.25pt\fi
    \ifnum\thesis@cellno=\@ne\ifthesis@leftflush\@tempdimb\z@\fi\fi
    \ifnum\thesis@cellno=\thesis@ncols\ifthesis@rightflush\@tempdimc\z@\fi\fi
  \fi}
\def\thesis@grey{\gdef\CT@row@color{\thesis@cellstep\CT@color{thesisstripe}}\global\thesis@curgreytrue}
\def\thesis@white{\gdef\CT@row@color{\thesis@cellstep}\global\thesis@curgreyfalse}
\ExplSyntaxOn
\cs_new_protected:Npn \thesis@edges #1
  {
    \str_set:Nn \l_tmpa_str {#1}
    \str_remove_all:Nn \l_tmpa_str { ~ }
    \str_if_eq:eeTF { \str_head:N \l_tmpa_str } { @ }
      { \global\thesis@leftflushtrue } { \global\thesis@leftflushfalse }
    \str_if_eq:eeTF { \str_range:Nnn \l_tmpa_str { -3 } { -1 } } { \tl_to_str:n { @{} } }
      { \global\thesis@rightflushtrue } { \global\thesis@rightflushfalse }
  }
\ExplSyntaxOff
\def\thesis@array[#1]#2{%
  \ifnum\thesis@lvl<\z@\thesis@edges{#2}\else\ifnum\currentgrouplevel<\thesis@lvl\thesis@edges{#2}\fi\fi
  \thesis@o@array[#1]{#2}}
\def\thesis@colour{\ifthesisstripes\ifodd\thesis@row\thesis@white\else\thesis@grey\fi\else\thesis@white\fi}
\def\thesis@startdata{\global\thesis@datatrue\global\thesis@head\m@ne\global\thesis@row\@ne
  \global\thesis@five\z@\thesis@colour}
\def\thesis@sethead{\global\thesis@lvl\currentgrouplevel\global\thesis@firstrulefalse
  \global\thesis@ncols\z@\global\thesis@cellno\z@
  \global\thesis@datafalse\global\thesis@row\z@\global\thesis@five\z@\thesis@white
  \global\thesis@head=\thesis@headrows\relax\ifnum\thesis@head=\z@\thesis@startdata\fi}

% ---- at every row end (\CT@everycr); only the outer table's rows count
\def\thesis@close{\ifnum0=`{\fi}}
\def\thesis@cr{\noalign{\ifnum0=`}\fi\thesis@crstep}
\def\thesis@crstep{\thesis@wantsepfalse
  \ifnum\currentgrouplevel=\thesis@lvl
    \ifthesis@curgrey\global\thesis@lastgreytrue\else\global\thesis@lastgreyfalse\fi
    \global\thesis@lastrowtrue
    \ifnum\thesis@cellno>\thesis@ncols\global\thesis@ncols\thesis@cellno\fi
    \global\thesis@cellno\z@\global\thesis@mc\z@
    \ifthesis@data
      \global\advance\thesis@row\@ne\thesis@colour
      \ifthesisfiverules
        \global\advance\thesis@five\@ne
        \ifnum\thesis@five<\thesisfiveevery\relax\else\global\thesis@five\z@\thesis@wantseptrue\fi
      \fi
    \else
      \thesis@white
      \ifnum\thesis@head>\z@
        \global\advance\thesis@head\m@ne
        \ifnum\thesis@head=\z@\thesis@startdata\fi
      \fi
    \fi
  \fi
  \ifthesis@wantsep\expandafter\thesis@peek\else\expandafter\thesis@close\fi}
% a thin rule after every \thesisfiveevery body rows, unless a rule or the end comes next
\def\thesis@peek{\futurenonspacelet\thesis@next\thesis@decide}
\def\thesis@nosepif#1{\ifx\thesis@next#1\thesis@oksepfalse\fi}
\def\thesis@decide{\thesis@okseptrue
  \thesis@nosepif\bottomrule\thesis@nosepif\midrule\thesis@nosepif\toprule\thesis@nosepif\cmidrule
  \thesis@nosepif\hline\thesis@nosepif\cline\thesis@nosepif\addlinespace\thesis@nosepif\specialrule
  \thesis@nosepif\end\thesis@nosepif\noalign\thesis@nosepif\crcr
  \ifthesis@oksep{\color{thesisfiverule}\hrule\@height\cmidrulewidth}\fi
  \thesis@close}

% ---- rules: #1 = 0 \cmidrule, 1 \toprule, 2 \midrule, 3 \bottomrule, 4 \hline
\def\thesis@fill{{\color{thesisstripe}\hrule\@height\aboverulesep}\vskip-\aboverulesep}
\def\thesis@rule#1{\noalign{\thesis@rulestep#1}}
\def\thesis@rulestep#1{%
  \ifnum#1=\@ne\thesis@sethead\else
  \ifnum#1=4
    \ifthesis@firstrule\thesis@sethead\else
    \ifnum\currentgrouplevel=\thesis@lvl\thesis@startdata\else\thesis@sethead\fi\fi
  \else
  \ifnum\currentgrouplevel=\thesis@lvl
    \ifthesis@lastrow\ifthesis@lastgrey\thesis@fill\fi\fi
    \ifnum#1=\tw@\thesis@startdata\fi
    \ifnum#1=\thr@@\global\thesis@datafalse\thesis@white\fi
  \fi\fi\fi
  \global\thesis@lastrowfalse}

% ---- at the start of each tabular / tabularx / tabular* (not nested ones): reset, taller rows
\def\thesis@tabbegin{%
  \ifthesis@active
    \ifnum\thesis@lvl<\z@\thesis@tabreset\else
    \ifnum\currentgrouplevel<\thesis@lvl\thesis@tabreset\fi\fi
  \fi}
\def\thesis@tabreset{\global\thesis@lvl\m@ne\global\thesis@firstruletrue\global\thesis@datafalse
  \global\thesis@row\z@\global\thesis@five\z@\global\thesis@head\m@ne\global\thesis@lastrowfalse
  \global\thesis@curgreyfalse
  \edef\arraystretch{\fpeval{\arraystretch*\thesisrowstretch}}}
\AddToHook{env/tabular/begin}{\thesis@tabbegin}
\AddToHook{env/tabular*/begin}{\thesis@tabbegin}
\AddToHook{env/tabularx/begin}{\thesis@tabbegin}

% ---- the style itself
\newcommand\thesistablestyle[1][-1]{%
  \def\thesis@headrows{#1}%
  \ifthesis@active\else
    \thesis@activetrue\global\thesis@lvl\m@ne
    \thesis@enginefalse
    \ifthesisstripes\thesis@enginetrue\fi
    \ifthesisfiverules\thesis@enginetrue\fi
    \ifthesis@engine
      \CT@everycr{\the\everycr\thesis@cr}%
      \let\thesis@o@toprule\toprule\def\toprule{\thesis@rule1\thesis@o@toprule}%
      \let\thesis@o@midrule\midrule\def\midrule{\thesis@rule2\thesis@o@midrule}%
      \let\thesis@o@bottomrule\bottomrule\def\bottomrule{\thesis@rule3\thesis@o@bottomrule}%
      \let\thesis@o@cmidrule\cmidrule\def\cmidrule{\thesis@rule0\thesis@o@cmidrule}%
      \let\thesis@o@hline\hline\def\hline{\thesis@rule4\thesis@o@hline}%
      \ifdefined\@@array\let\thesis@o@array\@@array\let\@@array\thesis@array\fi
      \ifdefined\mcell@@@tabular\let\thesis@o@mcell\mcell@@@tabular\let\mcell@@@tabular\thesis@mcell\fi
      \let\thesis@o@docolor\CT@@do@color\let\CT@@do@color\thesis@docolor
    \fi
  \fi}
\newcommand\thesisnostripes{%
  \ifthesis@active
    \thesis@activefalse
    \ifthesis@engine
      \CT@everycr{\noalign{\global\let\CT@row@color\relax}\the\everycr}%
      \let\toprule\thesis@o@toprule\let\midrule\thesis@o@midrule\let\bottomrule\thesis@o@bottomrule
      \let\cmidrule\thesis@o@cmidrule\let\hline\thesis@o@hline
      \ifdefined\thesis@o@array\let\@@array\thesis@o@array\fi
      \ifdefined\thesis@o@mcell\let\mcell@@@tabular\thesis@o@mcell\fi
      \let\CT@@do@color\thesis@o@docolor
      \thesis@enginefalse
    \fi
  \fi}
\def\thesis@auto{\ifthesistablesauto\thesistablestyle\fi}
\AddToHook{env/table/begin}{\thesis@auto}
\AddToHook{env/table*/begin}{\thesis@auto}
\AddToHook{env/sidewaystable/begin}{\thesis@auto}
\AddToHook{env/sidewaystable*/begin}{\thesis@auto}
\makeatother

%
% What it does, in every table / table* / sidewaystable / sidewaystable* float:
%   * zebra striping of the body rows (header rows stay white), in a very light grey that prints in greyscale;
%   * slightly taller rows (\arraystretch times \thesisrowstretch, on top of any \arraystretch the table sets);
%   * booktabs rules keep their spacing; where a rule follows a shaded row, the gap above the rule is shaded
%     too, so there is no white band.
% The body starts after the first \midrule (booktabs) or the second \hline (\hline tables). Striping restarts
% after every \midrule, so a group row after a \midrule (e.g. "\multicolumn{8}{l}{\textit{PropNEAT}}") is
% white and its rows alternate below it. A table whose header ends without a \midrule says how many header
% rows it has: \thesistablestyle[1]. A table with no top rule (\toprule or \hline) is not striped.
%
% Switches (in the preamble for the whole thesis, or inside one table for that table):
%   \thesisstripestrue / \thesisstripesfalse        zebra on (default) / off
%   \thesisfiverulestrue / \thesisfiverulesfalse    thin grey rule every \thesisfiveevery body rows (default off);
%                                                   the alternative to stripes: \thesisstripesfalse\thesisfiverulestrue
%   \thesistablesautotrue / \thesistablesautofalse  apply to every table float automatically (default on); when
%                                                   off, only tables that call \thesistablestyle are styled
%   \renewcommand\thesisrowstretch{1.1}             row-height factor (1 = unchanged)
%   \renewcommand\thesisfiveevery{5}                rows between the thin rules
%   \definecolor{thesisstripe}{gray}{0.93}          stripe colour (and thesisfiverule for the thin rules)
% Per table:
%   \thesistablestyle[<n>]   apply the style (generated tables call it first; idempotent); <n> = number of
%                            header rows after \toprule when the header does not end with a \midrule
%   \thesisnostripes         opt out: this table keeps its own look (no stripes, no thin rules, no extra
%                            height). Use it for tables that colour their own cells or rows, or use \multirow
%                            (the next row's shading would cover the \multirow text).
% Nested tabulars (\makecell, a tabular inside a cell) are left alone and do not disturb the outer striping.
\makeatletter
\@ifpackageloaded{colortbl}{}{\RequirePackage{colortbl}}
\@ifpackageloaded{booktabs}{}{\RequirePackage{booktabs}}
\@ifundefined{fpeval}{\RequirePackage{xfp}}{}

\newif\ifthesisstripes   \thesisstripestrue
\newif\ifthesisfiverules \thesisfiverulesfalse
\newif\ifthesistablesauto \thesistablesautotrue
\definecolor{thesisstripe}{gray}{0.93}
\definecolor{thesisfiverule}{gray}{0.6}
\newcommand\thesisrowstretch{1.1}
\newcommand\thesisfiveevery{5}

\newif\ifthesis@active   % inside a styled float
\newif\ifthesis@engine   % stripes or thin rules are being drawn
\newif\ifthesis@data     % past the header
\newif\ifthesis@firstrule
\newif\ifthesis@lastrow  % the last item was a row (not a rule)
\newif\ifthesis@curgrey  % the current row is shaded
\newif\ifthesis@lastgrey % the row just ended was shaded
\newif\ifthesis@wantsep
\newif\ifthesis@oksep
\newcount\thesis@row     % body row number since the last \midrule (1 = first)
\newcount\thesis@five
\newcount\thesis@head    % header rows still to come (-1: wait for a \midrule)
\newcount\thesis@lvl     % group level of the outer table's \noalign (-1: unknown)
\thesis@lvl=\m@ne
\def\thesis@headrows{-1}

% ---- row colours (colortbl's per-row colour, as xcolor's \rowcolors sets it). colortbl extends each
% cell's colour by \tabcolsep on both sides; at an @{} edge that would stick out past the rules, so the
% first / last cell of a row drops that overhang when the column spec starts / ends with @{}.
\newif\ifthesis@leftflush
\newif\ifthesis@rightflush
\newcount\thesis@cellno
\newcount\thesis@ncols
% colortbl stretches a cell's colour with 1fill against the cell's own 1fill of alignment glue; a
% \makecell adds its own \hfill glue (counted in \thesis@mc), which would take part of the stretch and
% leave the right of the cell white. The colour's stretch is raised to match (\thesis@k).
\newcount\thesis@mc
\def\thesis@k{1}
\def\thesis@mcell#1#2\@nil#3{\thesis@o@mcell{#1}#2\@nil{#3}%
  \ifx\mcell@left\hfill\global\advance\thesis@mc\@ne\fi
  \ifx\mcell@right\hfill\global\advance\thesis@mc\@ne\fi}
\def\thesis@docolor{\global\let\CT@do@color\relax
  \@tempdima\wd\z@\advance\@tempdima\@tempdimb\advance\@tempdima\@tempdimc
  \kern-\@tempdimb\leaders\vrule\hskip\@tempdima\@plus\thesis@k fill
  \kern-\@tempdimc\hskip-\wd\z@\@plus-\thesis@k fill\relax}
\def\thesis@cellstep{%
  \def\thesis@k{1}%
  \ifnum\currentgrouplevel=\thesis@lvl
    \edef\thesis@k{\the\numexpr\@ne+\thesis@mc\relax}\global\thesis@mc\z@
    \global\advance\thesis@cellno\@ne
    % overlap the next cell a little, so screen renderers show no hairline gaps between cells
    \ifnum\thesis@cellno<\thesis@ncols\advance\@tempdimc.25pt\fi
    \ifnum\thesis@cellno=\@ne\ifthesis@leftflush\@tempdimb\z@\fi\fi
    \ifnum\thesis@cellno=\thesis@ncols\ifthesis@rightflush\@tempdimc\z@\fi\fi
  \fi}
\def\thesis@grey{\gdef\CT@row@color{\thesis@cellstep\CT@color{thesisstripe}}\global\thesis@curgreytrue}
\def\thesis@white{\gdef\CT@row@color{\thesis@cellstep}\global\thesis@curgreyfalse}
\ExplSyntaxOn
\cs_new_protected:Npn \thesis@edges #1
  {
    \str_set:Nn \l_tmpa_str {#1}
    \str_remove_all:Nn \l_tmpa_str { ~ }
    \str_if_eq:eeTF { \str_head:N \l_tmpa_str } { @ }
      { \global\thesis@leftflushtrue } { \global\thesis@leftflushfalse }
    \str_if_eq:eeTF { \str_range:Nnn \l_tmpa_str { -3 } { -1 } } { \tl_to_str:n { @{} } }
      { \global\thesis@rightflushtrue } { \global\thesis@rightflushfalse }
  }
\ExplSyntaxOff
\def\thesis@array[#1]#2{%
  \ifnum\thesis@lvl<\z@\thesis@edges{#2}\else\ifnum\currentgrouplevel<\thesis@lvl\thesis@edges{#2}\fi\fi
  \thesis@o@array[#1]{#2}}
\def\thesis@colour{\ifthesisstripes\ifodd\thesis@row\thesis@white\else\thesis@grey\fi\else\thesis@white\fi}
\def\thesis@startdata{\global\thesis@datatrue\global\thesis@head\m@ne\global\thesis@row\@ne
  \global\thesis@five\z@\thesis@colour}
\def\thesis@sethead{\global\thesis@lvl\currentgrouplevel\global\thesis@firstrulefalse
  \global\thesis@ncols\z@\global\thesis@cellno\z@
  \global\thesis@datafalse\global\thesis@row\z@\global\thesis@five\z@\thesis@white
  \global\thesis@head=\thesis@headrows\relax\ifnum\thesis@head=\z@\thesis@startdata\fi}

% ---- at every row end (\CT@everycr); only the outer table's rows count
\def\thesis@close{\ifnum0=`{\fi}}
\def\thesis@cr{\noalign{\ifnum0=`}\fi\thesis@crstep}
\def\thesis@crstep{\thesis@wantsepfalse
  \ifnum\currentgrouplevel=\thesis@lvl
    \ifthesis@curgrey\global\thesis@lastgreytrue\else\global\thesis@lastgreyfalse\fi
    \global\thesis@lastrowtrue
    \ifnum\thesis@cellno>\thesis@ncols\global\thesis@ncols\thesis@cellno\fi
    \global\thesis@cellno\z@\global\thesis@mc\z@
    \ifthesis@data
      \global\advance\thesis@row\@ne\thesis@colour
      \ifthesisfiverules
        \global\advance\thesis@five\@ne
        \ifnum\thesis@five<\thesisfiveevery\relax\else\global\thesis@five\z@\thesis@wantseptrue\fi
      \fi
    \else
      \thesis@white
      \ifnum\thesis@head>\z@
        \global\advance\thesis@head\m@ne
        \ifnum\thesis@head=\z@\thesis@startdata\fi
      \fi
    \fi
  \fi
  \ifthesis@wantsep\expandafter\thesis@peek\else\expandafter\thesis@close\fi}
% a thin rule after every \thesisfiveevery body rows, unless a rule or the end comes next
\def\thesis@peek{\futurenonspacelet\thesis@next\thesis@decide}
\def\thesis@nosepif#1{\ifx\thesis@next#1\thesis@oksepfalse\fi}
\def\thesis@decide{\thesis@okseptrue
  \thesis@nosepif\bottomrule\thesis@nosepif\midrule\thesis@nosepif\toprule\thesis@nosepif\cmidrule
  \thesis@nosepif\hline\thesis@nosepif\cline\thesis@nosepif\addlinespace\thesis@nosepif\specialrule
  \thesis@nosepif\end\thesis@nosepif\noalign\thesis@nosepif\crcr
  \ifthesis@oksep{\color{thesisfiverule}\hrule\@height\cmidrulewidth}\fi
  \thesis@close}

% ---- rules: #1 = 0 \cmidrule, 1 \toprule, 2 \midrule, 3 \bottomrule, 4 \hline
\def\thesis@fill{{\color{thesisstripe}\hrule\@height\aboverulesep}\vskip-\aboverulesep}
\def\thesis@rule#1{\noalign{\thesis@rulestep#1}}
\def\thesis@rulestep#1{%
  \ifnum#1=\@ne\thesis@sethead\else
  \ifnum#1=4
    \ifthesis@firstrule\thesis@sethead\else
    \ifnum\currentgrouplevel=\thesis@lvl\thesis@startdata\else\thesis@sethead\fi\fi
  \else
  \ifnum\currentgrouplevel=\thesis@lvl
    \ifthesis@lastrow\ifthesis@lastgrey\thesis@fill\fi\fi
    \ifnum#1=\tw@\thesis@startdata\fi
    \ifnum#1=\thr@@\global\thesis@datafalse\thesis@white\fi
  \fi\fi\fi
  \global\thesis@lastrowfalse}

% ---- at the start of each tabular / tabularx / tabular* (not nested ones): reset, taller rows
\def\thesis@tabbegin{%
  \ifthesis@active
    \ifnum\thesis@lvl<\z@\thesis@tabreset\else
    \ifnum\currentgrouplevel<\thesis@lvl\thesis@tabreset\fi\fi
  \fi}
\def\thesis@tabreset{\global\thesis@lvl\m@ne\global\thesis@firstruletrue\global\thesis@datafalse
  \global\thesis@row\z@\global\thesis@five\z@\global\thesis@head\m@ne\global\thesis@lastrowfalse
  \global\thesis@curgreyfalse
  \edef\arraystretch{\fpeval{\arraystretch*\thesisrowstretch}}}
\AddToHook{env/tabular/begin}{\thesis@tabbegin}
\AddToHook{env/tabular*/begin}{\thesis@tabbegin}
\AddToHook{env/tabularx/begin}{\thesis@tabbegin}

% ---- the style itself
\newcommand\thesistablestyle[1][-1]{%
  \def\thesis@headrows{#1}%
  \ifthesis@active\else
    \thesis@activetrue\global\thesis@lvl\m@ne
    \thesis@enginefalse
    \ifthesisstripes\thesis@enginetrue\fi
    \ifthesisfiverules\thesis@enginetrue\fi
    \ifthesis@engine
      \CT@everycr{\the\everycr\thesis@cr}%
      \let\thesis@o@toprule\toprule\def\toprule{\thesis@rule1\thesis@o@toprule}%
      \let\thesis@o@midrule\midrule\def\midrule{\thesis@rule2\thesis@o@midrule}%
      \let\thesis@o@bottomrule\bottomrule\def\bottomrule{\thesis@rule3\thesis@o@bottomrule}%
      \let\thesis@o@cmidrule\cmidrule\def\cmidrule{\thesis@rule0\thesis@o@cmidrule}%
      \let\thesis@o@hline\hline\def\hline{\thesis@rule4\thesis@o@hline}%
      \ifdefined\@@array\let\thesis@o@array\@@array\let\@@array\thesis@array\fi
      \ifdefined\mcell@@@tabular\let\thesis@o@mcell\mcell@@@tabular\let\mcell@@@tabular\thesis@mcell\fi
      \let\thesis@o@docolor\CT@@do@color\let\CT@@do@color\thesis@docolor
    \fi
  \fi}
\newcommand\thesisnostripes{%
  \ifthesis@active
    \thesis@activefalse
    \ifthesis@engine
      \CT@everycr{\noalign{\global\let\CT@row@color\relax}\the\everycr}%
      \let\toprule\thesis@o@toprule\let\midrule\thesis@o@midrule\let\bottomrule\thesis@o@bottomrule
      \let\cmidrule\thesis@o@cmidrule\let\hline\thesis@o@hline
      \ifdefined\thesis@o@array\let\@@array\thesis@o@array\fi
      \ifdefined\thesis@o@mcell\let\mcell@@@tabular\thesis@o@mcell\fi
      \let\CT@@do@color\thesis@o@docolor
      \thesis@enginefalse
    \fi
  \fi}
\def\thesis@auto{\ifthesistablesauto\thesistablestyle\fi}
\AddToHook{env/table/begin}{\thesis@auto}
\AddToHook{env/table*/begin}{\thesis@auto}
\AddToHook{env/sidewaystable/begin}{\thesis@auto}
\AddToHook{env/sidewaystable*/begin}{\thesis@auto}
\makeatother
